% An article, from TeXCraft's own templates (CC0-1.0: use it as you like).
\documentclass[11pt]{article}
\usepackage[margin=2.5cm]{geometry}
\usepackage{amsmath}
\usepackage{graphicx}
\usepackage{booktabs}
\usepackage[hidelinks]{hyperref}

\title{The Title of the Article}
\author{First Author \and Second Author}
\date{\today}

\begin{document}
\maketitle

\begin{abstract}
One paragraph on the question, what was done, and what was found.
\end{abstract}

\section{Introduction}
\label{sec:intro}
Say why the question matters and what this article adds. Section~\ref{sec:method}
describes the method, and Section~\ref{sec:results} the results.

\section{Method}
\label{sec:method}
A displayed equation, numbered so that the text can refer to it:
\begin{equation}
  E = mc^2 .
  \label{eq:energy}
\end{equation}
Equation~\eqref{eq:energy} is the best known in physics.

\section{Results}
\label{sec:results}
\begin{table}[ht]
  \centering
  \caption{A table with rules from booktabs.}
  \label{tab:results}
  \begin{tabular}{lrr}
    \toprule
    Case & Value & Error \\
    \midrule
    First  & 1.00 & 0.05 \\
    Second & 2.50 & 0.10 \\
    \bottomrule
  \end{tabular}
\end{table}
Table~\ref{tab:results} gives the values. Drag an image from the project's files onto this text
to insert it.

\section{Conclusion}
What was shown, and what remains open.

\begin{thebibliography}{9}
\bibitem{knuth} D.~E. Knuth. \emph{The \TeX book}. Addison-Wesley, 1984.
\end{thebibliography}

\end{document}
